\chapter{العالم الكمومي: الفوتونات ومستويات الطاقة}\label{ch:g12:quantum-world}

يقدر مصباح حراري \omterm{def:g10:light-spectra:wavelength}{تحت أحمر} أن يغمر جلدك طوال المساء ولا يترك فيه أثرًا؛ وعشر
دقائق من ضوء شمس جبلي رقيق تترك حرقًا. والمصباح يوصّل \omterm{def:g10:energy-conservation:energy}{طاقة} أكبر
بكثير --- لكن \omterm{def:g10:energy-conservation:energy}{الطاقة} ليست ما يترك أثرًا في الجلد. فالضوء، كما يبيّن هذا الفصل،
يصل حبيبات، ولا تجري الكيمياء إلا بحبيبة تضرب \omterm{def:g10:inertia:force}{بقوة} كافية:
فحبيبة فوق بنفسجية واحدة تقصف رابطة لا يفعل بها مليار حبيبة تحت حمراء إلا التسخين.
اتبع الحبيبات ينفتح قرن كامل: معادن تنزّ إلكترونات، وذرات
تغنّي في خطوط طيفية، وإلكترونات تتداخل كالأمواج.

\section{الضوء يصل حبيبات: الفوتون}

\begin{definition}[الفوتون]\label{def:g12:quantum-world:photon}
الضوء ذو \omterm{def:g12:oscillators-and-time:oscillator}{التواتر} $\nu$ لا يتبادل \omterm{def:g10:energy-conservation:energy}{الطاقة} مع المادة إلا بحبيبات كاملة
تسمّى \emph{الفوتونات}\index{فوتون}، \omterm{def:g10:energy-conservation:energy}{طاقة} كل منها
\[
E = h\nu = \frac{hc}{\lambda}, \qquad h = \qty{6.63e-34}{J.s},
\]
حيث $h$ هو \emph{ثابت بلانك}\index{ثابت بلانك} (بلانك،
1900؛ أينشتاين، 1905)، و$\lambda$ \omterm{def:g12:mechanical-waves:wavelength}{طول الموجة} في الخلاء. وشدة
الحزمة تقيس \emph{عدد} الفوتونات في \omterm{def:g10:orders-of-magnitude:unit}{الثانية}؛ أما \omterm{def:g10:energy-conservation:energy}{طاقة} كل حبيبة
فيضبطها \omterm{def:g12:oscillators-and-time:oscillator}{التواتر} وحده.
\end{definition}

\begin{notation}[الإلكترون فولط]\label{not:g12:quantum-world:ev}
تستعمل الفيزياء الذرية \emph{الإلكترون فولط}\index{إلكترون فولط}،
$\qty{1}{eV} = \qty{1.60e-19}{J}$: وهو \omterm{def:g10:energy-conservation:energy}{الطاقة} التي تكسبها \omterm{def:g11:fundamental-interactions:elementary}{شحنة أولية} واحدة
عابرةً \qty{1}{V} (\cref{ch:g11:electric-gravitational-fields}). ومن المفيد:
$E \,(\unit{eV}) \approx 1240/\lambda\,(\unit{nm})$.
\end{notation}

\begin{omfigure}
\begin{tikzpicture}[font=\small, x=1.48cm]
  \draw[-Stealth] (7.55,0) -- (-0.55,0);
  \draw[-Stealth, omDef] (4.9,-0.85) -- (2.3,-0.85)
    node[midway, below] {الطاقة لكل فوتون تتزايد};
  \foreach \x/\l/\e/\n in {0/{\qty{1}{pm}}/{\qty{1.2}{MeV}}/{$\gamma$},
      1.2/{\qty{0.1}{nm}}/{\qty{12}{keV}}/{X},
      2.4/{\qty{100}{nm}}/{\qty{12}{eV}}/{UV},
      3.6/{\qty{550}{nm}}/{\qty{2.3}{eV}}/{visible},
      4.8/{\qty{10}{\micro m}}/{\qty{0.12}{eV}}/{IR},
      6.0/{\qty{1}{cm}}/{\qty{0.12}{meV}}/{micro},
      7.2/{\qty{3}{m}}/{\qty{0.4}{\micro eV}}/{radio}}{
    \draw (\x,0.07) -- (\x,-0.07);
    \node[above, omExo] at (\x,0.42) {\n};
    \node[above, font=\scriptsize] at (\x,0.10) {\l};
    \node[below, font=\scriptsize, omDef] at (\x,-0.10) {\e};}
\end{tikzpicture}

{\small سُلّم واحد للعائلة \omterm{prop:g10:signals-and-waves:em}{الكهرمغناطيسية}
(\cref{ch:g12:light-as-wave}): \omterm{def:g12:mechanical-waves:wavelength}{طول الموجة} (في الأعلى)، \omterm{def:g10:energy-conservation:energy}{والطاقة} لكل \omterm{def:g12:quantum-world:photon}{فوتون}
$hc/\lambda$ (بالأزرق)، من حبيبات الراديو ذات الميكرو \omterm{def:g12:nuclear-energy:units}{إلكترون فولط} إلى
حبيبات \omterm{def:g11:nucleus-radioactivity:abg}{غاما} \unit{MeV} ذات \cref{ch:g11:nucleus-radioactivity} --- ومع ذلك يصبّ مؤشّر
قدره \qty{1.0}{mW} عددًا قدره $\num{3e15}$ حبيبة في \omterm{def:g10:orders-of-magnitude:unit}{الثانية}: فضوء
كل يوم يجري في سلاسة الرمل المصبوب.}
\end{omfigure}

\section{المفعول الكهرضوئي}

سلّط ضوءًا على صفيحة معدنية نظيفة في الخلاء، فتقفز إلكترونات إلى الخارج --- وهو
\emph{المفعول الكهرضوئي}\index{مفعول كهرضوئي}. وقواعده كسرت الفيزياء الكلاسيكية.

\begin{proposition}[ما تبيّنه التجربة]\label{prop:g12:quantum-world:facts}
لكل معدن \emph{تواتر عتبي}\index{تواتر عتبي} $\nu_0$:
\begin{enumerate}
  \item دون $\nu_0$، لا يغادر أي إلكترون، مهما تكن شدة الضوء؛
  \item وفوق $\nu_0$، يبدأ الإصدار بلا \omterm{def:g12:mechanical-waves:delay}{تأخير}، مهما يكن الضوء خافتًا؛
  \item والشدة ترفع \emph{عدد} الإلكترونات في \omterm{def:g10:orders-of-magnitude:unit}{الثانية}، لا
        طاقتها الحركية العظمى؛ \omterm{def:g12:oscillators-and-time:oscillator}{والتواتر} يرفع $E_{k,\max}$.
\end{enumerate}
ولا شيء من ذلك يوافق \omterm{def:g10:signals-and-waves:wave}{موجة} مستمرة، إذ ينبغي لأي لون أن يشحن ---
على مدى أشهر في ضوء خافت (\cref{exo:g12:quantum-world:14}) --- وينبغي لشدتها
أن تضبط \omterm{def:g10:energy-conservation:energy}{طاقة} الإلكترونات: فالتنبؤات الثلاثة كلها تفشل.
\end{proposition}
\admitted

\begin{theorem}[علاقة أينشتاين الكهرضوئية]\label{thm:g12:quantum-world:einstein}
\index{علاقة كهرضوئية}يحتاج الإلكترون على الأقل \emph{شغل
الخروج}\index{شغل الخروج} $W_0$ في المعدن \omterm{rem:g12:work-and-energy:escape}{ليفلت} من سطحه.
ويُمتصّ \omterm{def:g12:quantum-world:photon}{الفوتون} الواحد كاملًا في إلكترون واحد، ومن ثم
\[
h\nu = W_0 + E_{k,\max}, \qquad \text{ومنه } \nu_0 = \frac{W_0}{h}.
\]
\end{theorem}

\begin{proof}
محاسبة \omterm{def:g10:energy-conservation:energy}{الطاقة}: $h\nu$ داخلة، و$W_0$ منفقة على الخروج، والباقي حركي ---
أي $h\nu - W_0$ على الأكثر، وأقل من ذلك للإلكترونات التي تبدأ أعمق. ودون $\nu_0$
لا تقدر حبيبة واحدة أن تدفع رسم الخروج، مهما يكن عددها.
\end{proof}

\begin{definition}[كمون الإيقاف]\label{def:g12:quantum-world:stopping}
لقياس $E_{k,\max}$، يكبح \omterm{def:g11:circuits-and-power:voltage}{توتر} عكسي الإلكترونات
(\cref{ch:g11:electric-gravitational-fields})؛ و\emph{كمون
الإيقاف}\index{كمون الإيقاف} $U_s$ هو أصغر \omterm{def:g11:circuits-and-power:voltage}{توتر} يلغي
التيار: إذ يفشل أسرع الإلكترونات في العبور بالكاد، ومن ثم $E_{k,\max} = e\,U_s$.
\end{definition}

\begin{omfigure}
\begin{tikzpicture}
\begin{axis}[omaxis, width=8.8cm, height=5.4cm,
    xlabel={$\nu$ ($10^{14}\,\unit{Hz}$)}, ylabel={$E_{k,\max}$ (\unit{eV})},
    xmin=0, xmax=13.4, ymin=-3.1, ymax=3.2,
    xtick={2,4,6,8,10,12}, ytick={-2,-1,1,2,3}]
  \addplot[omExo, dashed, domain=0:5.55, samples=2] {0.414*x-2.3};
  \addplot[omThm, thick, domain=5.55:12.4, samples=2] {0.414*x-2.3};
  \addplot[only marks, mark=*, mark size=1.5pt, omDef] coordinates
    {(7,0.60) (8,1.01) (9,1.43) (10,1.84) (11,2.26)};
  \node[below, font=\small] at (axis cs:5.55,-0.08) {$\nu_0$};
  \node[right, font=\small, omExo] at (axis cs:0.15,-2.55) {$-W_0$};
  \node[font=\small, omThm, rotate=22] at (axis cs:10.3,1.35) {الميل $= h$};
\end{axis}
\end{tikzpicture}

{\small قياس $E_{k,\max} = eU_s$ بدلالة \omterm{def:g12:oscillators-and-time:oscillator}{التواتر} (في الصوديوم): أي
مستقيم ميله $h$، يقطع المحور عند $\nu_0$؛ ويكشف، ممدَّدًا
(بالمتقطع)، عن $-W_0$ عند $\nu = 0$.}
\end{omfigure}

\begin{method}[قياس ثابت بلانك]\label{met:g12:quantum-world:measure-h}
\begin{enumerate}
  \item أضِئ خلية كهرضوئية عبر مرشّحات معلومة أطوال الموجات؛
        واحسب كل $\nu = c/\lambda$؛ وسجّل كل كمون إيقاف.
  \item ارسم $E_{k,\max} = eU_s$ بدلالة $\nu$: فأينشتاين يعد
        بالمستقيم $h\nu - W_0$.
  \item اقرأ الميل: فهو \emph{يساوي} $h$؛ والقاطعان يعطيان $\nu_0$ و$-W_0$.
\end{enumerate}
\end{method}

\section{مستويات الطاقة: لماذا تصدر الذرات خطوطًا}

\begin{definition}[مستويات الطاقة]\label{def:g12:quantum-world:levels}
\omterm{def:g10:energy-conservation:energy}{طاقة} الإلكترون المرتبط في \omterm{def:g11:fundamental-interactions:ladder}{ذرة}
\emph{مكمَّمة}\index{تكميم}: أي محصورة في سُلّم منفصل من
\emph{مستويات الطاقة}\index{مستوى طاقة} $E_1 < E_2 < \dots < 0$، معدودةً
سالبةً انطلاقًا من الإلكترون الحر. وأخفض درجة هي \emph{الحالة
الأساسية}\index{حالة أساسية}، وسواها \emph{الحالات
المثارة}\index{حالة مثارة}؛ وبلوغ $0$ هو
\emph{التأيّن}\index{تأيّن}.
\end{definition}

\begin{proposition}[سُلّم الهيدروجين]\label{prop:g12:quantum-world:hydrogen}
\omterm{def:g12:quantum-world:levels}{مستويات الطاقة} في \omterm{def:g11:fundamental-interactions:ladder}{ذرة} الهيدروجين هي
\[
E_n = -\frac{\qty{13.6}{eV}}{n^2}, \qquad n = 1, 2, 3, \dots
\]
ومن ثم $E_1 = \qty{-13.6}{eV}$، $E_2 = \qty{-3.40}{eV}$، $E_3 =
\qty{-1.51}{eV}$، $E_4 = \qty{-0.85}{eV}$، متزاحمةً نحو $0$.
\end{proposition}
\admitted

\begin{remark}[من أين يأتي السُّلّم]\label{rem:g12:quantum-world:ladder-origin}
أما لماذا هذه الأعداد فميكانيك كمومي، محفوظ لمجلدات الجامعة ---
وكذلك صيغة دي بروي: \omterm{def:g12:quantum-world:debroglie}{فموجة مادة} الإلكترون
(\cref{def:g12:quantum-world:debroglie} أدناه) يجب أن تنغلق على نفسها حول
\omterm{def:g11:fundamental-interactions:ladder}{النواة}، ولا يوافق ذلك إلا أطوال \omterm{def:g10:signals-and-waves:wave}{موجة} معينة --- فللذرة، كوتر الغيتار،
مجموعة منفصلة من النغمات.
\end{remark}

\begin{theorem}[فوتون انتقال]\label{thm:g12:quantum-world:transition}
لا تغيّر \omterm{def:g11:fundamental-interactions:ladder}{الذرة} مستواها إلا بإصدار \omterm{def:g12:quantum-world:photon}{فوتون} واحد أو امتصاصه يحمل
الفرق بالضبط: فالهبوط من المستوى $p$ إلى مستوى أخفض $n$ يصدر
$h\nu = E_p - E_n$؛ والصعود من $n$ إلى $p$ يمتصّ \omterm{def:g10:energy-conservation:energy}{الطاقة} نفسها.
\end{theorem}

\begin{proof}
\omterm{thm:g10:energy-conservation:conservation}{انحفاظ الطاقة}، حبيبةً حبيبة: \omterm{def:g12:quantum-world:photon}{فالفوتون} هو العملة الوحيدة التي
يقبلها الضوء، ودفاتر \omterm{def:g11:fundamental-interactions:ladder}{الذرة} يجب أن تتوازن.
\end{proof}

\begin{remark}[تفسير الأطياف الخطية]\label{rem:g12:quantum-world:lines}
هذا هو الجواب الموعود به في \cref{ch:g10:light-spectra}: فالغاز الحارّ يصدر
تواترات سُلّمه الخاص $(E_p - E_n)/h$ لا غير --- أي رمزًا شريطيًّا من
الخطوط اللامعة --- والغاز البارد يسرق تلك التواترات بالضبط من \omterm{def:g10:light-spectra:white}{الضوء
الأبيض}، فيطبع خطوط امتصاص داكنة مطابقة. ولكل عنصر سُلّمه
الخاص، ومن ثم رمزه الشريطي الخاص --- وهكذا نقرأ ضوء النجوم.
\end{remark}

\begin{example}[خطوط الهيدروجين المرئية]\label{ex:g12:quantum-world:balmer}
تقع الانتقالات المنتهية عند $n = 2$ في المرئي. ومع $E \,(\unit{eV}) =
1240/\lambda\,(\unit{nm})$: $3 \to 2$: \qty{1.89}{eV}، \qty{656}{nm} (أحمر)؛
$4 \to 2$: \qty{2.55}{eV}، \qty{486}{nm} (أزرق مخضرّ)؛ $5 \to 2$:
\qty{2.86}{eV}، \qty{434}{nm} (بنفسجي). أما الانتقالات إلى $n = 1$ فتتجاوز
\qty{10}{eV}: أي فوق بنفسجية كلها.
\end{example}

\begin{omfigure}
\begin{tikzpicture}[font=\small, y=1cm]
  \foreach \y/\n/\e in {0/1/-13.6, 4.50/2/-3.40, 5.33/3/-1.51}{
    \draw[thick] (0,\y) -- (4.2,\y);
    \node[left] at (0,\y) {$n=\n$};
    \node[right, font=\scriptsize] at (4.2,\y) {\qty{\e}{eV}};}
  \draw[thick] (0,5.63) -- (4.2,5.63) node[left,at start,yshift=-2.5pt] {$n=4$};
  \draw[thick] (0,5.76) -- (4.2,5.76) node[left,at start,yshift=3.5pt] {$n=5$};
  \draw[dashed] (0,6.0) -- (4.2,6.0) node[left, at start] {$\infty$}
    node[right, font=\scriptsize] {\qty{0}{eV} (متأيّنة)};
  \foreach \x/\ys/\col/\l in {1.0/5.33/omThm/656, 1.9/5.63/omDef/486,
      2.8/5.76/violet/434} \draw[-Stealth, \col, thick] (\x,\ys) --
    (\x,4.50) node[midway, left, font=\scriptsize] {\qty{\l}{nm}};
  \draw[-Stealth, omExo, thick] (3.7,4.50) -- (3.7,0)
    node[midway, right, font=\scriptsize] {\qty{122}{nm} (فوق بنفسجي)};
\end{tikzpicture}

{\small سُلّم الهيدروجين، بالمقياس: تنتهي الخطوط المرئية كلها عند
$n = 2$؛ والهبوط إلى \omterm{def:g12:quantum-world:levels}{الحالة الأساسية} \omterm{def:g10:light-spectra:wavelength}{فوق بنفسجي} عميق.}
\end{omfigure}

\begin{remark}[الليزر، في نفَس واحد]\label{rem:g12:quantum-world:laser}
\omterm{def:g12:quantum-world:photon}{الفوتون} ذو \omterm{def:g10:energy-conservation:energy}{الطاقة} $E_p - E_n$ بالضبط، إذا مرّ \omterm{def:g11:fundamental-interactions:ladder}{بذرة} \emph{مثارة}، قدر أن يحرّضها
على إصدار \omterm{def:g12:quantum-world:photon}{فوتون} مطابق --- \omterm{def:g10:energy-conservation:energy}{بالطاقة} نفسها، والاتجاه نفسه، وعلى الخطى. وأبقِ
عدد الذرات في الأعلى أكبر منه في الأسفل، يصر \omterm{def:g12:quantum-world:photon}{فوتون} واحد انهيارًا من النسخ:
أي \emph{الإصدار المستحثّ}\index{إصدار مستحثّ} --- \omterm{def:g11:color-light-sources:lamps}{والليزر} ليس أكثر من ذلك بكثير.
\end{remark}

\section{أمواج المادة}

\begin{proposition}[الإلكترونات تتداخل]\label{prop:g12:quantum-world:interference}
أرسل الإلكترونات واحدًا واحدًا عبر شقّين
(\cref{ch:g12:light-as-wave}): فيهبط كلٌّ منها نقطة حادّة واحدة --- أي جسيمًا.
لكن مع تراكم الآلاف، ترسم النقاط أهداب تداخل:
فكل إلكترون، وحده، يمرّ موجةً عبر \emph{كلا} الشقين. والبلورات
تحيد حزم الإلكترونات أيضًا (دافيسون وجيرمر، 1927).
\end{proposition}
\admitted

\begin{omfigure}
\begin{tikzpicture}[font=\small]
  \foreach \o in {0,4.0,8.0} \draw[omExo] (\o,0) rectangle +(3.4,2.2);
  \foreach \x/\y in {0.62/1.71, 1.68/0.42, 2.31/1.95, 1.05/1.22,
      2.98/0.66, 1.74/1.58, 0.38/0.35, 2.40/0.98, 1.62/1.02, 0.97/1.86,
      2.95/1.60, 1.71/0.20, 1.12/0.61, 2.33/0.31}
    \fill (\x,\y) circle (0.8pt);
  \foreach \c/\f in {0.40/610, 1.05/410, 1.70/530, 2.35/350, 3.00/470}
    \foreach \y in {0.25,0.62,...,2.00}
      \fill ({4.0+\c+0.14*sin(\f*\y)},\y) circle (0.8pt);
  \foreach \x/\y in {4.72/1.42, 5.38/0.52, 6.06/1.12, 6.66/1.72,
      5.44/1.88, 7.28/0.78} \fill (\x,\y) circle (0.8pt);
  \foreach \c/\f in {0.40/610, 1.05/410, 1.70/530, 2.35/350, 3.00/470}
    \foreach \y in {0.12,0.27,...,2.10}
      \fill ({8.0+\c+0.11*sin(\f*\y)},\y) circle (0.8pt);
  \foreach \c/\f in {0.40/260, 1.05/340, 1.70/205, 2.35/290, 3.00/385}
    \foreach \y in {0.20,0.50,...,2.05}
      \fill ({8.0+\c+0.06*sin(\f*\y+40)},\y) circle (0.8pt);
  \node[below] at (1.7,-0.05) {$\sim 30$ إلكترونًا};
  \node[below] at (5.7,-0.05) {$\sim 200$}; \node[below] at (9.7,-0.05) {$\sim 2000$};
\end{tikzpicture}

{\small إلكترون واحد، ونقطة واحدة؛ وآلاف، وأهداب. وأين تهبط النقطة التالية
أمر عشوائي --- كتفكك \omterm{def:g11:fundamental-interactions:ladder}{نواة} واحدة (\cref{ch:g11:nucleus-radioactivity})
--- لكن \omterm{def:g10:signals-and-waves:wave}{الموجة} تثبّت \emph{الاحتمال}: فالهدب اللامع مرجَّح، والداكن
لا يكاد يحدث. أي صدفة فردية يحكمها موج مضبوط: وهو قلب الفيزياء الكمومية.}
\end{omfigure}

\begin{definition}[طول موجة دي بروي]\label{def:g12:quantum-world:debroglie}
كل جسيم كمية حركته $p = mv$ يسافر \emph{موجة
مادة}\index{موجة مادة} ذات \emph{طول موجة دي بروي}\index{طول موجة دي
بروي}
\[
\lambda = \frac{h}{p} = \frac{h}{mv}.
\]
\end{definition}
\admitted

\begin{example}[لماذا لا تحيد أنت أبدًا]\label{ex:g12:quantum-world:never}
الإلكترون المسرَّع عبر \qty{100}{V} له $E_k = \qty{1.60e-17}{J}$،
$p = \sqrt{2m_e E_k} = \qty{5.4e-24}{kg.m/s}$، ومن ثم $\lambda \approx
\qty{0.12}{nm}$ --- أي تباعد الذرات في بلورة، وهي تعمل من ثم
عمل شقَّيه. أما كرة تنس كتلتها \qty{58}{g} \omterm{def:g12:kinematics-2d:velocity}{بسرعة} \qty{25}{m/s}: فيكون $\lambda =
6.63\times10^{-34}/1.45 \approx \qty{4.6e-34}{m}$، أي أصغر $10^{19}$ مرة
من \omterm{def:g11:fundamental-interactions:ladder}{نواة} --- فلن يظهرها أي شقّ تُهدِّب أبدًا. و$h$ صغير جدًّا
بحيث لا تطفو الحبيبية الكمومية إلا على مقياس \omterm{def:g11:fundamental-interactions:ladder}{الذرة}.
\end{example}

\begin{example}[المجهر الإلكتروني]\label{ex:g12:quantum-world:microscope}
يمنع \omterm{def:g12:light-as-wave:diffraction}{الحيود} تمييز تفصيل أصغر بكثير من \omterm{def:g12:mechanical-waves:wavelength}{طول الموجة} المستعمل
(\cref{ch:g12:light-as-wave}): فالضوء المرئي يتوقف قرب \qty{0.5}{\micro m}.
وللإلكترونات عند \qty{10}{kV} المقدار $\lambda \approx \qty{12}{pm}$ --- أي أقصر عشرات
الآلاف من المرات: فالمجاهر الإلكترونية تصوّر الفيروسات، وآلات
الخلية، حتى الذرات المفردة. وقد سلّمت \omterm{def:g12:quantum-world:debroglie}{موجة المادة} الفيزياءَ عينًا جديدة.
\end{example}

\section{تمارين}

\begin{exercise}[$\star$]\label{exo:g12:quantum-world:1}
مؤشّر \omterm{def:g11:color-light-sources:lamps}{ليزر} أخضر يصدر \qty{1.0}{mW} عند \qty{530}{nm}. احسب \omterm{def:g10:energy-conservation:energy}{طاقة}
\omterm{def:g12:quantum-world:photon}{فوتون} واحد \omterm{def:g11:work-of-force:work}{بالجول} وبالإلكترون \omterm{def:g11:circuits-and-power:voltage}{فولط}، ثم \omterm{def:g12:quantum-world:photon}{الفوتونات} في \omterm{def:g10:orders-of-magnitude:unit}{الثانية}.
\end{exercise}

\begin{exercise}[$\star$]\label{exo:g12:quantum-world:2}
محطة FM تبثّ عند \qty{100}{MHz}. احسب \omterm{def:g10:energy-conservation:energy}{طاقة} فوتونها \omterm{def:g10:orders-of-magnitude:unit}{بوحدة}
\unit{eV} وقارنها \omterm{def:g12:quantum-world:photon}{بفوتون} مرئي: فلماذا لا تُكتشف حبيبية الراديو؟
\end{exercise}

\begin{exercise}[$\star$]\label{exo:g12:quantum-world:3}
للسيزيوم $W_0 = \qty{1.9}{eV}$. احسب تواتره العتبي \omterm{def:g10:light-spectra:wavelength}{وطول موجته}
العتبي. وهل يقذف \omterm{def:g11:color-light-sources:lamps}{ليزر} أحمر عند \qty{633}{nm} إلكترونات من السيزيوم ---
وإن كان كذلك، فبأي \omterm{def:g11:mechanical-energy:kinetic}{طاقة حركية} عظمى؟
\end{exercise}

\begin{exercise}[$\star$]\label{exo:g12:quantum-world:4}
باستعمال $E_n = -13.6/n^2$ \unit{eV}، احسب \omterm{def:g10:energy-conservation:energy}{طاقة} \omterm{def:g12:quantum-world:photon}{الفوتون} المصدَر
في انتقال الهيدروجين $3 \to 2$ \omterm{def:g10:light-spectra:wavelength}{وطول موجته}. وما لونه؟
\end{exercise}

\begin{exercise}[$\star$]\label{exo:g12:quantum-world:5}
من مخطط $E_{k,\max}$ بدلالة $\nu$ في الدرس، اقرأ \omterm{prop:g12:quantum-world:facts}{التواتر
العتبي} للصوديوم وشغل خروجه (\omterm{def:g10:orders-of-magnitude:unit}{بوحدة} \unit{eV})؛ وأيّ ثابت هو الميل؟
\end{exercise}

\begin{exercise}[$\star\star$]\label{exo:g12:quantum-world:6}
ضوء \omterm{def:g10:light-spectra:wavelength}{طول موجته} \qty{400}{nm} يسقط على البوتاسيوم ($W_0 = \qty{2.3}{eV}$).
احسب \omterm{def:g10:energy-conservation:energy}{طاقة} \omterm{def:g12:quantum-world:photon}{الفوتون}، \omterm{def:g11:mechanical-energy:kinetic}{والطاقة الحركية} العظمى للإلكترونات
(\omterm{def:g10:orders-of-magnitude:unit}{بوحدة} \unit{eV} وبالجول)، \omterm{def:g12:quantum-world:stopping}{وكمون الإيقاف}.
\end{exercise}

\begin{exercise}[$\star\star$]\label{exo:g12:quantum-world:7}
خلية كهرضوئية مضاءة فوق العتبة. فماذا يحدث (أ) للتيار، (ب) و
\omterm{def:g12:quantum-world:stopping}{لكمون الإيقاف}، حين تتضاعف الشدة عند \omterm{def:g12:oscillators-and-time:oscillator}{تواتر} ثابت؟ وحين
يُرفع \omterm{def:g12:oscillators-and-time:oscillator}{التواتر}؟ وأيّ واقعة تناقض الصورة الموجية، ولماذا؟
\end{exercise}

\begin{exercise}[$\star\star$]\label{exo:g12:quantum-world:8}
أيّ \omterm{def:g10:energy-conservation:energy}{طاقة} \omterm{def:g12:quantum-world:photon}{فوتون} دنيا تؤيّن \omterm{def:g11:fundamental-interactions:ladder}{ذرة} هيدروجين من حالتها الأساسية؟ ومن
$n = 2$؟ احسب طولَي \omterm{def:g10:signals-and-waves:wave}{الموجة} العتبيين وسمِّ مجاليهما الطيفيين.
\end{exercise}

\begin{exercise}[$\star\star$]\label{exo:g12:quantum-world:9}
\omterm{def:g10:light-spectra:white}{ضوء أبيض} يعبر وعاءً من هيدروجين بارد (وكل الذرات في \omterm{def:g12:quantum-world:levels}{الحالة الأساسية}).
فسّر لماذا لا يظهر \omterm{def:g12:sound-acoustics:timbre}{الطيف} النافذ \emph{أي} خط داكن في
المرئي، بينما تظهر خطوط امتصاص بالمر في أطياف النجوم الحارّة.
\end{exercise}

\begin{exercise}[$\star\star$]\label{exo:g12:quantum-world:10}
احسب \omterm{def:g12:quantum-world:debroglie}{طول موجة دي بروي} في الحالتين: (أ) إلكترون مسرَّع عبر
\qty{100}{V}، (ب) ذبابة كتلتها \qty{0.10}{g} \omterm{def:g12:kinematics-2d:velocity}{بسرعة} \qty{1.0}{m/s}. وقارن كلًّا منهما
بحجم \omterm{def:g11:fundamental-interactions:ladder}{ذرة} ($\sim \qty{e-10}{m}$)؛ واستنتج.
\end{exercise}

\begin{exercise}[$\star\star$]\label{exo:g12:quantum-world:11}
تستجيب العين المتكيفة مع الظلام لنحو $90$ \omterm{def:g12:quantum-world:photon}{فوتون} من ضوء \qty{505}{nm}.
احسب تلك \omterm{def:g10:energy-conservation:energy}{الطاقة}، وقارنها \omterm{def:g11:mechanical-energy:kinetic}{بالطاقة الحركية} لبعوضة كتلتها \qty{2.5}{mg}
تطير \omterm{def:g12:kinematics-2d:velocity}{بسرعة} \qty{0.50}{m/s}، وفسّر النسبة.
\end{exercise}

\begin{exercise}[$\star\star\star$]\label{exo:g12:quantum-world:12}
خلية كهرضوئية تعطي $U_s = \qty{1.19}{V}$ عند \qty{405}{nm} و$U_s =
\qty{0.40}{V}$ عند \qty{546}{nm}. ومن هاتين النقطتين قِس $h$، ثم
\omterm{thm:g12:quantum-world:einstein}{شغل الخروج} \omterm{def:g10:orders-of-magnitude:unit}{بوحدة} \unit{eV}. وأيّ معدن من هذا الفصل هو المهبط؟
\end{exercise}

\begin{exercise}[$\star\star\star$]\label{exo:g12:quantum-world:13}
مصباح هيدروجين يظهر خطين عند \qty{486}{nm} و\qty{434}{nm}. حوّل
كلًّا منهما إلى \omterm{def:g10:energy-conservation:energy}{طاقة} \omterm{def:g12:quantum-world:photon}{فوتون} وعيّن الانتقالين بمطابقة
فروق المستويات $E_n = -13.6/n^2$ \unit{eV}.
\end{exercise}

\begin{exercise}[$\star\star\star$]\label{exo:g12:quantum-world:14}
تقدير \omterm{def:g12:mechanical-waves:delay}{التأخير} الكلاسيكي: ضوء خافت شدته \qty{1.0e-6}{W/m^2}
يسقط على البوتاسيوم ($W_0 = \qty{2.3}{eV}$). فإذا جمعت \omterm{def:g11:fundamental-interactions:ladder}{ذرة} \omterm{def:g10:energy-conservation:energy}{الطاقة}
العابرة مقطعها العرضي وحدها (ونصف قطره \qty{1.0e-10}{m})، فكم من الزمن
تحتاج لتجمع $W_0$؟ قارن \omterm{def:g12:mechanical-waves:delay}{بالتأخير} المرصود (دون \qty{e-9}{s}).
\end{exercise}

\begin{exercise}[$\star\star\star$]\label{exo:g12:quantum-world:15}
في مجهر إلكتروني تُسرَّع الإلكترونات عبر \qty{5.0}{kV}.
احسب \omterm{def:g12:quantum-world:debroglie}{طول موجة دي بروي} لها، وقارنه بالضوء الأخضر (\qty{550}{nm})،
وفسّر \omterm{def:g12:light-as-wave:diffraction}{بالحيود} (\cref{ch:g12:light-as-wave}) لماذا يشتري ذلك تمييزًا.
\end{exercise}

\section{مسألة: مختبر فيزياء اللوح الشمسي}

\begin{problem}[{مسألة نهاية الأسبوع --- مختبر فيزياء اللوح الشمسي:
حبيبات قدرها $10^{21}$ يشربها سطح في كل \omterm{def:g10:orders-of-magnitude:unit}{ثانية}، وخلية كهرضوئية تقيس
\omterm{def:g12:quantum-world:photon}{ثابت بلانك}، ومصباح هيدروجين يُقرأ خطًّا خطًّا --- نهاية أسبوع واحدة، و$h$ واحد}]
\label{pb:g12:quantum-world:1}
يجهّز صفّك اللوح الشمسي على سطح المدرسة نهايةَ أسبوع.
والمعطيات: الشدة الشمسية الساقطة عند الظهر \qty{1.0e3}{W/m^2}؛ \omterm{def:g12:mechanical-waves:wavelength}{وطول الموجة} الشمسي المتوسط
\qty{550}{nm}؛ ومساحة اللوح \qty{1.0}{m^2}، \omterm{def:g10:energy-conservation:efficiency}{ومردوده} $20\%$؛ واليوم
يوصّل ما يعادل \qty{5.0}{h} من الشمس الكاملة.

\medskip
\noindent\textbf{الجزء I --- محاسبة \omterm{def:g12:quantum-world:photon}{الفوتونات}.}
\begin{enumerate}
  \item احسب \omterm{def:g10:energy-conservation:energy}{طاقة} \omterm{def:g12:quantum-world:photon}{فوتون} شمسي متوسط واحد، \omterm{def:g11:work-of-force:work}{بالجول} وبالإلكترون \omterm{def:g11:circuits-and-power:voltage}{فولط}.
  \item استنتج عدد \omterm{def:g12:quantum-world:photon}{الفوتونات} التي تضرب اللوح في \omterm{def:g10:orders-of-magnitude:unit}{الثانية} عند الظهر.
  \item احسب \omterm{prop:g11:circuits-and-power:power}{الاستطاعة الكهربائية} للوح عند الظهر، ثم \omterm{def:g10:energy-conservation:energy}{الطاقة}
        المنتَجة على مدى اليوم، \omterm{def:g10:orders-of-magnitude:unit}{بوحدة} \unit{kW.h}.
  \item ويبدو تيار اللوح سلسًا تمامًا. وباستعمال السؤال 2،
        فسّر لماذا يكون الوصول حبيبةً حبيبة غير مرئي.
  \item ويحمل ضوء الشمس \omterm{def:g10:energy-conservation:energy}{طاقة} تحت حمراء أكبر من الفوق بنفسجية، ومع ذلك لا يحرق
        الجلد إلا \omterm{def:g10:light-spectra:wavelength}{الفوق بنفسجي} (فالروابط $\sim \qtyrange{3}{4}{eV}$). فسّر بالحبيبات.
\end{enumerate}

\medskip
\noindent\textbf{الجزء II --- قياس $h$ بالخلية الكهرضوئية في العدّة.}
للخلية الكهرضوئية في العدّة مهبط مجهول. وتعطي المرشّحات:
\begin{center}\small
\begin{tabular}{lcccc}
$\lambda$ (\unit{nm}) & 365 & 405 & 436 & 546\\
$U_s$ (\unit{V}) & 1.50 & 1.16 & 0.94 & 0.37
\end{tabular}
\end{center}
\begin{enumerate}[resume]
  \item ماذا يقيس \omterm{def:g12:quantum-world:stopping}{كمون الإيقاف}؟ اربط $U_s$ و$\nu$ و$W_0$ و$h$.
  \item حوّل أطوال الموجات الأربعة إلى تواترات (\omterm{def:g10:orders-of-magnitude:unit}{بوحدة} $10^{14}\,\unit{Hz}$).
  \item ولماذا ينبغي أن يعطي رسم $eU_s$ بدلالة $\nu$ مستقيمًا؟
        تحقق من اصطفاف النقاط.
  \item من النقطتين الطرفيتين، احسب الميل: أي $h$ المقيس عندك.
  \item استنتج شغل خروج المهبط (\unit{eV})، وتواتره
        العتبي، \omterm{def:g10:light-spectra:wavelength}{وطول موجته} العتبي.
  \item وأيّ جزء من ضوء الشمس يقدر أن يقذف إلكترونات من هذا
        المهبط --- وأيّ نصف من \omterm{def:g12:sound-acoustics:timbre}{الطيف} عديم النفع له؟
  \item قارن $h$ عندك بالقيمة \qty{6.63e-34}{J.s}: فما الخطأ النسبي بالمئة؟
\end{enumerate}

\medskip
\noindent\textbf{الجزء III --- مصباح الهيدروجين للمعايرة.}
يُتحقق من سُلّم أطوال الموجات في العدّة في مقابل مصباح هيدروجين،
$E_n = -13.6/n^2$ \unit{eV}.
\begin{enumerate}[resume]
  \item احسب $E_2$ و$E_3$ و$E_4$ و$E_5$.
  \item احسب طاقات فوتونات الانتقالات $3 \to 2$ و$4 \to 2$ و$5 \to 2$.
  \item استنتج أطوال الموجات الثلاثة وأعطِ لون كل خط.
  \item ولماذا يصدر المصباح \emph{خطوطًا} لا قوس قزح
        مستمرًّا؟ جملة واحدة، تسمّي الفكرة المفتاح.
  \item وهل يقدر هيدروجين المصباح أن يمتصّ ضوءًا عند \qty{500}{nm}؟ برّر بالسُّلّم.
\end{enumerate}

\medskip
\noindent\textbf{الجزء IV --- فحص الخلايا: \omterm{def:g10:signals-and-waves:wave}{موجة} الإلكترون.}
الشقوق المجهرية في خلية تقع دون متناول \qty{0.5}{\micro m} في المجاهر
الضوئية بكثير؛ والمجهر الإلكتروني الماسح في المختبر يعمل عند \qty{10}{kV}.
\begin{enumerate}[resume]
  \item احسب \omterm{def:g12:kinematics-2d:velocity}{سرعة} الإلكترونات؛ وتحقق من بقائها دون $0.25\,c$ (بالصيغة الكلاسيكية).
  \item احسب كمية حركتها \omterm{def:g12:quantum-world:debroglie}{وطول موجة دي بروي} لها.
  \item والتمييز يتناسب مع \omterm{def:g12:mechanical-waves:wavelength}{طول الموجة}: احسب المكسب على الضوء الأخضر
        (\qty{550}{nm})؛ واختم برقمَي نهاية الأسبوع --- أي $h$ عندك
        وهذا المكسب.
\end{enumerate}
\end{problem}
